% !TeX root = ../sections/02_exercises.tex

\begin{exercise}[情報数学 A 中間試験 2]
	周期 \(2\pi\) をもち，区間 \([-\pi,\pi)\) 上で
	\[
	f(x)=
	\begin{cases}
		0
		& \left(|x|\leq \dfrac{\pi}{3}\right),\\[6pt]
		\dfrac{\pi}{2}
		& \left(\dfrac{\pi}{3}<|x|\leq \dfrac{2\pi}{3}\right),\\[6pt]
		\pi
		& \left(\dfrac{2\pi}{3}<|x|\leq \pi\right)
	\end{cases}
	\]
	で与えられる周期関数について，次の問いに答えよ．
	
	\begin{enumerate}
		\item \(y=f(x)\) のグラフをかけ．
		\item \(f(x)\) のフーリエ係数を求めよ．
		\item \(f(x)\) のフーリエ級数の第 \(n\) 部分和を
		\[
		S_n(x)=\frac{a_0}{2}
		+\sum_{k=1}^{n}
		\left(
		a_k\cos kx+b_k\sin kx
		\right)
		\]
		で表す．gnuplot を用いて
		\[
		y=f(x),\qquad y=S_7(x),\qquad y=S_{31}(x)
		\]
		のグラフを一つの図にプロットした \texttt{.png} ファイル
		\texttt{graph.png} を作成せよ．
	\end{enumerate}
	
	ただし，次のようにすること．
	
	\begin{itemize}
		\item グラフの色は任意であるが，見やすい色を指定することが望ましい．
		\item \(y=f(x)\) のグラフがうまくプロットできない場合は省略せよ
		（その場合，減点はする）．
		\item \(x\) 軸の範囲はデフォルトの設定である \([-10,10]\) のまま，
		\(y\) 軸の範囲は \([-0.5,4]\) とせよ．
		\item \(y=f(x)\)，\(y=S_7(x)\)，\(y=S_{31}(x)\) のグラフの凡例は
		それぞれ \(\mathrm{f(x)}\)，\(\mathrm{S7(x)}\)，\(\mathrm{S31(x)}\) とせよ．
		\item グラフをプロットする際はスクリプトファイル \texttt{plot.txt} を作成し，
		それを \texttt{load} することでプロットせよ．ただし，
		「作成した図を \texttt{.png} ファイルに書き出す命令文」は
		\texttt{plot.txt} に含めないこと．
		\item 作成した二つのファイル
		\[
		\texttt{graph.png},\qquad \texttt{plot.txt}
		\]
		を学務情報システムの提出物にアップロードせよ．
	\end{itemize}
\end{exercise}
\begin{background}
	この問題では，周期 \(2\pi\) の段階関数
	\[
	f(x)=
	\begin{cases}
		0
		& \left(|x|\leq \dfrac{\pi}{3}\right),\\[6pt]
		\dfrac{\pi}{2}
		& \left(\dfrac{\pi}{3}<|x|\leq \dfrac{2\pi}{3}\right),\\[6pt]
		\pi
		& \left(\dfrac{2\pi}{3}<|x|\leq \pi\right)
	\end{cases}
	\]
	を扱う。
	
	この関数は \(|x|\) によって定義されているため，偶関数である。
	すなわち，
	\[
	f(-x)=f(x)
	\]
	が成り立つ。
	
	したがって，フーリエ級数を求める際には正弦係数がすべて消える。
	つまり，
	\[
	b_n=0
	\]
	である。
	
	フーリエ係数の一般公式は
	\[
	a_0=\frac{1}{\pi}\int_{-\pi}^{\pi}f(x)\,dx,
	\]
	\[
	a_n=\frac{1}{\pi}\int_{-\pi}^{\pi}f(x)\cos nx\,dx,
	\]
	\[
	b_n=\frac{1}{\pi}\int_{-\pi}^{\pi}f(x)\sin nx\,dx
	\]
	である。
	
	ただし，今回は偶関数なので，
	\[
	a_0=\frac{2}{\pi}\int_0^\pi f(x)\,dx,
	\]
	\[
	a_n=\frac{2}{\pi}\int_0^\pi f(x)\cos nx\,dx,
	\]
	\[
	b_n=0
	\]
	として計算できる。
	
	また，この関数は区分的に定数である。
	したがって積分は，関数の値が変化する点
	\[
	0,\quad \frac{\pi}{3},\quad \frac{2\pi}{3},\quad \pi
	\]
	で区間を分割して行う。
	
	グラフについては，中央で \(0\)，その外側で \(\pi/2\)，さらに外側で \(\pi\)
	となる左右対称な階段状のグラフである。
	
	gnuplot では，このような段階関数を条件分岐で表す。
	さらに，\([-10,10]\) のように複数周期を含む範囲で描画するためには，
	任意の \(x\) を基本区間 \([-\pi,\pi)\) に戻してから関数値を決める必要がある。
\end{background}
\begin{strategy}
	この問題は，三つの小問を一つの流れとして考えるとよい。
	
	まず，グラフを描く。
	関数は \(|x|\) の大きさで値が決まるので，グラフは \(y\)-軸対称である。
	
	区間ごとの値は次の通りである。
	\[
	|x|\leq \frac{\pi}{3}
	\quad\Longrightarrow\quad
	f(x)=0,
	\]
	\[
	\frac{\pi}{3}<|x|\leq \frac{2\pi}{3}
	\quad\Longrightarrow\quad
	f(x)=\frac{\pi}{2},
	\]
	\[
	\frac{2\pi}{3}<|x|\leq \pi
	\quad\Longrightarrow\quad
	f(x)=\pi.
	\]
	
	したがって，\(-\pi\) から \(\pi\) までの範囲では，
	中央が低く，外側に行くほど高くなる階段状のグラフを描けばよい。
	
	次に，フーリエ係数を求める。
	偶関数であることから，まず
	\[
	b_n=0
	\]
	と分かる。
	
	次に，
	\[
	a_0=\frac{2}{\pi}\int_0^\pi f(x)\,dx
	\]
	を計算する。
	
	関数値ごとに積分区間を分けると，
	\[
	a_0
	=
	\frac{2}{\pi}
	\left(
	\int_0^{\pi/3}0\,dx
	+
	\int_{\pi/3}^{2\pi/3}\frac{\pi}{2}\,dx
	+
	\int_{2\pi/3}^{\pi}\pi\,dx
	\right)
	\]
	となる。
	
	次に，
	\[
	a_n=\frac{2}{\pi}\int_0^\pi f(x)\cos nx\,dx
	\]
	を同様に区間分割して，
	\[
	a_n
	=
	\frac{2}{\pi}
	\left(
	\int_{\pi/3}^{2\pi/3}\frac{\pi}{2}\cos nx\,dx
	+
	\int_{2\pi/3}^{\pi}\pi\cos nx\,dx
	\right)
	\]
	とする。
	
	ここで
	\[
	\int \cos nx\,dx=\frac{\sin nx}{n}
	\]
	を使えばよい。
	
	計算結果をまとめると，
	\[
	a_0=\pi,
	\]
	\[
	a_n
	=
	-\frac{1}{n}
	\left(
	\sin\frac{n\pi}{3}
	+
	\sin\frac{2n\pi}{3}
	\right),
	\]
	\[
	b_n=0
	\]
	となる。
	
	したがって，第 \(N\) 部分和は
	\[
	S_N(x)
	=
	\frac{\pi}{2}
	+
	\sum_{n=1}^{N}
	a_n\cos nx
	\]
	である。
	
	gnuplot では，この \(a_n\) と \(S_N(x)\) をそのまま関数として定義する。
	重要なのは次の三つである。
	
	\begin{enumerate}
		\item 基本区間に戻す関数を定義する。
		\item 段階関数 \(f(x)\) を条件分岐で定義する。
		\item \(S_7(x)\) と \(S_{31}(x)\) を \texttt{sum} を使って定義する。
	\end{enumerate}
	
	基本区間に戻すには，たとえば
	\[
	p(x)=x-2\pi\left\lfloor\frac{x+\pi}{2\pi}\right\rfloor
	\]
	を使う。
	これにより，任意の \(x\) を \([-\pi,\pi)\) に戻すことができる。
	
	その上で，\(|p(x)|\) の値によって
	\[
	0,\quad \frac{\pi}{2},\quad \pi
	\]
	を返すようにすれば，周期関数としての \(f(x)\) が描ける。
	
	この問題の発想は，
	\[
	\text{偶関数}
	\Rightarrow
	b_n=0
	\Rightarrow
	a_n\text{ だけ区間分割で計算}
	\Rightarrow
	S_N(x)\text{ を gnuplot で描画}
	\]
	という流れである。
\end{strategy}